\chapter[The Islamic World's Connections]{How the Islamic World Linked the Old World}
\section{Three Things in Your Exercise Book}
Before turning the page, look beside you.

Perhaps your maths homework lies open. Pick up three things from it, with your eyes rather than your hands.

First: paper. Thin, light, cheap, disposable without regret. You probably never consider that writing material was once among humanity's greatest expenses. Parchment used entire animal skins; a thick book could consume a flock. Chinese bamboo slips required ox-carts; Egyptian papyrus grew brittle after repeated rolling. Cheap, durable paper was first made in China, then travelled an extraordinarily long road to your desk.

Second: numerals, such as 37 and 1208. We call them Arabic numerals, a name concealing an injustice: Indians invented them, not Arabs. Arabs did something else: put them on the road.

Third: the equation you cannot solve. Your teacher calls the subject algebra and its procedures algorithms. Both words trace to a Baghdad scholar twelve hundred years ago: one from his book's title, the other from the Latin form of his name.

Why do paper, numerals, and equations meet in your homework? They travelled one network, extending at its broadest from Atlantic Spain to Central Asian Samarkand and Indian Ocean ports. A language, books, and perpetually travelling people connected it, not roads and walls alone. We call it the Islamic world, but that can mislead: Muslims, Christians, Jews, and Indian scholars wove it, carrying not only religion but business, mathematics, and curiosity.

This chapter asks how it was woven, what it carried, and most importantly whether it merely moved other people's possessions or created new things.

\section{A Baghdad Bookshop around 830}
Enter a scene.

Time: a day in the first half of the ninth century, roughly China's late Tang. Place: Baghdad on the Tigris's east bank. Founded only decades earlier, it is already among the world's largest cities, with hundreds of thousands inside circular walls. Ships arriving from the Indian Ocean unload Chinese ceramics, Indian pepper, and East African ivory.

You are heading not to the docks but into an ink-scented shop. A warraq combines copyist and bookseller. Scribes sit cross-legged at low tables before paper---yes, Baghdad already uses it; more shortly---with sharpened reed pens and soot-based ink. Assistants sort and bind finished sheets behind them. Customers order poetry or astronomical tables as casually as you might order coffee in a bookshop.

A man around fifty frowns over two written sheets. His name is Hunayn ibn Ishaq. Look carefully: almost all this chapter's secrets meet in him.

First, he is Christian, not Muslim, belonging to the eastern community called Nestorian, condemned as heretical in Byzantium and rooted further east in Persia and Mesopotamia. Second, translation earns his living, chiefly rendering Galen's Greek medical works into Arabic, often through an initial Syriac version in his church's language, followed by an assistant's Arabic translation. A Greek doctor's sentence passes through two linguistic relays to reach Arab readers. Third, Muslim caliphs, ministers, and rich merchants compete to employ him. Later biography says one patron paid a manuscript's weight in gold. Possibly exaggerated, the story conveys a real point: knowledge of Greek was worth a gold mine here.

Outside, the city does similar work. Caliph al-Ma'mun established a scholarly institution in his palace, later called the House of Wisdom. Historians recount his dream of Aristotle and subsequent support for Greek philosophical translation. The dream's truth is unknown; the money was real. Bilingual and trilingual scholars arrived from Syria, Persia, and Central Asia: Christians translated medicine, descendants of Zoroastrians calculated astronomy, Persian descendants wrote philosophy, and Indian scholars explained remarkable numerals including a circular zero.

Remember the shop: a Christian, paid by a Muslim ruler, transfers Greek books through Syriac into Arabic, writing on Chinese-invented paper for imperial readers. Nobody here is culturally pure. That mixture wove an Old World network.

\section{The Real Question: How Was It Woven?}
Pull back far enough to see the whole map.

Westward, Córdoba's Islamic dynasty maintained a library of hundreds of thousands of volumes according to medieval accounts. Treat the number cautiously, though even a tenth would exceed any contemporary European monastery. Eastward, Samarkand's workshops made famous paper day and night. Southward, Red Sea vessels rode monsoons to India, guided by stars and navigation manuals. Between stood Cairo, Damascus, and Basra, cities of hundreds of thousands with religious schools, disputing theologians, and merchants working abacuses.

Riding from Córdoba to Samarkand took over half a year. Yet medical students at either end might use the same textbook in the same language, native to neither: one might speak a local Romance dialect, the other Persian. Their shared Arabic resembled scientists' shared English today.

Three real questions emerge:

\textbf{First: how was the network woven?} Arab armies conquered from Spain to Central Asia in the seventh and eighth centuries, but swords mark borders, not bookshops. Later Mongols conquered almost as much without weaving the same network. Was money the cause, religion, or something less visible?

\textbf{Second: were its goods imported or created?} A long-standing European account credits Arabs with preserving Greek knowledge like warehouse custodians, keeping books until Renaissance Europeans collected them. Courteous-sounding, but correct? Do custodians revise stock, identify its errors, or add entirely new products?

\textbf{Third: who decides centre and margin?} Baghdad was central, but Samarkand paper, Indian numerals, and Córdoban medicine defined the network from elsewhere. Arabic connected it while millions of Persians, Syrians, and Spaniards lowered their mother tongues' volume to participate. How many gained the world, and how many lost their voices?

None permits a quick answer. Take them one at a time.

\section[Rulers, Translators, and Traders]{Many Voices: Rulers, Translators, Merchants, and the Translated}
No single voice weaves a network. Listen to several.

\textbf{The caliph.} Scholarship was a shrewd investment. Having overthrown a dynasty, Abbasids needed legitimacy; sponsoring Greek and Persian translations proclaimed inheritance of ancient wisdom. Practical needs mattered too: doctors, tax and engineering calculations, precise calendars. To courts, learning was half prestige, half tool.

\textbf{The translator.} Hunayn's position was delicate: serving Muslim rulers, translating pagan Greeks, belonging to neither group. Such intermediaries' partial belonging enabled movement between languages. He translated over a hundred Galenic works and complained about finding good manuscripts, searching Syrian, Palestinian, and Mesopotamian monasteries for one book. Translation already combined detective work with physical exertion a thousand years ago.

\textbf{The Persian.} The network's largest minority possessed proud civilisational memories: Sasanian Ctesiphon had long gathered Greek and Indian learning. In the eighth century, Persian-descended Ibn al-Muqaffa translated a Pahlavi, or Old Persian, work derived from Indian fables into Arabic. Kalila and Dimna spread animal tales of political wisdom across Eurasia, its roots reaching India's Panchatantra. Hear the route: Indian stories in Persian clothing, then Arabic clothing, then languages worldwide. Even fables changed hands repeatedly.

\textbf{The merchant.} Scholars paid to travel; merchants travelled to profit, often aboard the same ship. A Cairo synagogue storeroom preserved tens of thousands of medieval letter fragments because Jewish custom stored rather than discarded written papers, accidentally creating the Cairo Geniza archive. Jewish merchants wrote from Cairo to partners on India's Malabar coast about schedules, pepper prices, debts, and family greetings, spelling their own language in Arabic letters. Not scholars, they nonetheless carried ledgers, bills, measures, and numeral systems across seas. Knowledge travelled with them.

\textbf{Religious scholars, and their disagreements.} Never imagine a single Islamic voice. Eleventh-century al-Ghazali's Incoherence of the Philosophers attacked Greek-inspired thinkers for overextending reason. Over a century later, Córdoba's Ibn Rushd answered with Incoherence of the Incoherence, defending philosophy while also serving as a judge and authority on Islamic law. Even the legitimacy of Greek philosophy divided devout scholars sharply.

\textbf{Finally, hear a voice you may oppose: the custodian theory.}

Give its case---that Islamic civilisation merely preserved Greece for Europe---the fullest hearing. Our recurring exercise strengthens an opposing argument before deciding whether to refute it.

Its evidence has force. First, without Arabic versions, much of Galen and parts of Aristotle might genuinely be lost: a fact, not flattery. Second, twelfth-century Europeans did recover Aristotle extensively through Arabic-to-Latin translation, often packaged with Arabic commentary. Third, the translation movement initially sought external learning; translations may numerically have exceeded originals.

Now examine the gaps.

First, custodians do not rewrite stock. Ibn Sina's Canon of Medicine reorganised, expanded, and corrected Galen rather than merely translating him, remaining a European university textbook into the seventeenth century. Students received upgraded Arabic medicine, not preserved Greek medicine alone. Second, custodians invent no new products, but al-Khwarizmi's algebra, Ibn al-Haytham's experimental optics, and al-Biruni's method for calculating Earth's radius were absent from Greece's warehouse. Third, the theory narrows suppliers: Indian mathematics and numerals, Persian administration, and Chinese papermaking also arrived. This was never a Greek speciality shop.

Preservation is true but insufficient. The network transported, unpacked, inspected, corrected, assembled, and created. We will establish that with evidence.

\section[Thinking Bridge: Map the Network]{Thinking Bridge: Draw Civilisation as a Network}
How can a topic spanning half the globe, seven or eight centuries, and many peoples avoid drowning us in detail? Our portable tool is \textbf{network mapping}.

Textbooks colour civilisations as bounded blocks: Chinese here, Islamic there, European elsewhere. This suggests boxes with fixed personalities. Our bookshop shows instead moving people and goods, without pure identities.

A network needs only three elements:

\textbf{First, nodes.} Draw cities as dots: Baghdad, Cairo, Córdoba, Samarkand, Basra, Damascus, and perhaps West African Timbuktu, where scholars copied, collected, and debated Arabic books for centuries at the Sahara's edge. Nodes vary in size; none is naturally central.

\textbf{Second, links.} Trade routes, pilgrimage paths carrying thousands annually to Mecca alongside books and news, and Indian Ocean monsoon routes. Summer and winter winds reverse, forcing merchants to wait in ports for months. Ledgers, algorithms, stories, and case histories change hands in teahouses and meeting places. Nature's rhythm becomes knowledge's rhythm.

\textbf{Third, travelling packages and three questions.} People, books, techniques, and words are packages. From which node to which? Who pays: courts, merchants, or churches? Most importantly, is the package opened and modified en route?

Practise on numerals. Starting in India, whose decimal notation included zero and whose seventh-century mathematicians discussed its arithmetic, they reached Baghdad, where al-Khwarizmi wrote about calculating with them, then North Africa and Spain, and finally European recommendation in Italian merchant Fibonacci's 1202 Liber Abaci. Courts partly paid, through House of Wisdom patronage; merchants partly paid themselves, needing arithmetic to avoid losses. Were the contents modified? Extensively: shapes changed during travel, leaving today's Indian numerals, Arabic numerals, and your 123 as differently dressed relatives.

The tool offers protection: \textbf{beware single-event origin stories}. You have probably heard that Chinese captives after the 751 Battle of Talas carried papermaking west. Tidy, but probably oversimplified. Paper may have reached Samarkand earlier, and techniques usually diffuse through prolonged contact, not one night's captives. Du Huan's account, shortly to be read, is real; a single battle changing technological history deserves a question mark. Existing links continuously carried packages without awaiting a great battle.

Networks also explain survival. Mongols captured Baghdad and killed its last caliph in 1258, removing a super-node. A civilisation-as-box would be smashed, but Cairo, Samarkand, and Córdoba continued operating as goods rerouted. Later Timur's grandson Ulugh Beg built a leading Samarkand observatory and remarkably accurate star tables. A hole in the net did not end the net.

Next time a chapter discusses a civilisation, redraw it as nodes, links, and modified packages. You will see what coloured blocks hide.

\section[Evidence: A Twelve-Hundred-Year-Old Equation]{Reading Evidence: An Equation Twelve Hundred Years Old}
Now touch primary material. Claims of new knowledge need evidence.

First: al-Khwarizmi's algebra book.

Working in early-ninth-century Baghdad's House of Wisdom, al-Khwarizmi came from Central Asian Khwarazm: another person moving from margin to centre. Around 820 he wrote a book whose title contained al-jabr, restoration, repairing something broken, even bone-setting in contemporary Arabic. Its Latin translation supplied Europeans with a subject name: al-jabr became algebra.

Near the opening, he gives a standard solution for the first equation type. Here is a paraphrase of its accepted surviving content, not a word-for-word translation:

\begin{quote}
A square plus ten roots equals thirty-nine dirhams, the silver coins of the time. Find the number.
\end{quote}
In familiar notation: x² + 10x = 39. His entirely verbal solution halves the roots' number: half ten is five. Square five to obtain twenty-five, add thirty-nine to get sixty-four, take its square root, eight, then subtract five. Three remains: the answer.

Check it: nine plus thirty is thirty-nine. Correct.

Pause over what happened. Instead of one isolated answer, he supplied a \textbf{procedure}: halve, square, add, take the square root, subtract. It applies throughout this equation class, whatever the numbers. This is an algorithm's simplest sense: steps independent of specific inputs that yield the answer. Three centuries later, a Latin translation of his book on Indian arithmetic began ``Algoritmi says,'' Latinising his name. Algoritmi became algorithm. Every recommendation and navigation route on your phone runs algorithms; the word preserves a Central Asian name in Baghdad.

The book's later examples concern inheritance, commerce, and land measurement, not mere puzzles. Complex inheritance shares divided estates among many relatives through fractions and equations. Courts and markets generated demand. He was not simply moving Greek books into Arabic: Greek equation-solving followed a very different geometrical style. He made a new tool for his world's problems.

Second: the eyes of a Tang captive.

In 751, Abbasid forces defeated Tang troops beside Central Asia's Talas River and captured soldiers and artisans. Du Huan spent over a decade in the western regions, returned to Guangzhou by merchant ship, and wrote the Record of Travels. The original is lost, but his relative Du You quoted fragments in the Comprehensive Institutions.

Du describes Dashi, the Arab Empire, including Tang artisans working as painters and silk weavers, naming some and their hometowns. He also describes customs: people worshipped neither king nor parents and elders, revering heaven alone, and on appointed days gathered in a great hall to hear a leader speak from above.

Why is this precious? It supplies eyes from the network's other end: an ordinary Tang person without a mission, captured into it, seeing its life. People themselves were travelling packages; captive workers carried enduring skills. But recall our caution: this does not establish papermaking's transmission at Talas. Du names painters and weavers, not papermakers. Accept what evidence says without inventing what it omits: the first discipline of primary-source reading.

\section{One Network, Three Explanations}
With a network, multilingual scholars, new mathematics, and a captive's testimony, we can address how it was woven and why it worked so vigorously. Three serious explanations concern why, not what: shared events, disputed causes.

\textbf{Explanation one: state investment (politics).}

The translation movement's peak coincided with Abbasid strength. Caliphs founded the House of Wisdom, paid generously, supported scholars, and sent collectors to Byzantium. Al-Ma'mun reportedly sought manuscripts in negotiations with its emperor. States needed doctors, astronomical tables, tax officials, and engineers, and financed knowledge accordingly. Funding records make this concrete. Yet it cannot explain other nodes: Córdoba and Cairo had golden ages under different rulers, and scholarship persisted in schools and private circles after courtly patronage changed. States lit fires, but fires burned beyond palace hearths.

\textbf{Explanation two: faith's impetus (religion).}

Islam values knowledge; believers understand scripture as inviting observation and reflection. Practice also creates technical needs. Facing Mecca thousands of kilometres away poses spherical geometry; five daily prayers require astronomical timing; Ramadan requires observing the new moon; inheritance requires shares. These stimulated astronomy, trigonometry, and algebra. Two gaps remain. Why so many non-Muslim contributors, such as Christian Hunayn or Jewish Maimonides, born in Córdoba, deceased in Cairo, author in Arabic of the Guide for the Perplexed? They need not face Mecca. And why internal disagreement such as al-Ghazali's and Ibn Rushd's? Faith may explain enthusiasm, not its disputes. Also beware the saying ``seek knowledge even in China.'' Scholars examining hadith transmission long questioned it, commonly treating it as fabricated or extremely weak. A knowledge-valuing civilisation must scrutinise forged praise of knowledge: excellent close-reading practice.

\textbf{Explanation three: commerce's routes.}

Across Spain and India, merchants needed accounts, currency exchange among dirhams and dinars of differing purity, monsoon schedules, and correspondence with distant partners. Geniza letters show literate, numerate, contract-observing networks demanding practical mathematics and geography while carrying books and ideas. This explains broadly why ports so often became scholarly cities. Yet business requires adequate arithmetic, while the network funded apparently impractical curiosity: Ibn al-Haytham's spherical-mirror reflections, al-Biruni's Earth's circumference and Indian cosmologies. Merchants needed none of these. Commerce alone cannot explain curiosity.

Each captures a real thread: courtly gold, faith's warp, commerce's weft. No single thread makes a net. This is not evasive compromise but a method: where several evidenced explanations each miss the whole, identify which portion each explains rather than choosing one winner.

\section[Further Challenge: Relay or Network?]{Further Challenge: A Relay Baton or a Network?}
Now take on the heaviest theoretical question, concerning \textbf{how we tell history itself}, not one historical detail.

A familiar framework presents civilisation as a relay. Greeks run first with philosophy, science, and democracy; Romans take over; medieval darkness nearly drops the baton; Arabs temporarily safeguard it; Renaissance Europe retrieves it and races to modernity. Call this the \textbf{relay-baton view of history}.

It reduces the Islamic world's role to handing over a baton, not running: standing beside the track holding something for tired runners.

Weigh the framework against this chapter's evidence.

First ask who made it. Renaissance Europeans sought ancestry bypassing contemporary church tradition to claim ancient Greece directly, blanking the millennium between ``us$\rightarrow$\allowbreak Greece.'' Nineteenth-century European expansion then wanted civilisation progressing one-way from west to east to explain Europe's position at the finish. This was a narrative for particular people's needs, not necessarily fabrication: it captured genuine reliance on Arabic Aristotle but mistook selected truth for the whole.

Recent historians of science identify three major faults.

First, there is no unchanged baton. Aristotle entering Arabic was annotated line by line, reclassified, and combined with Persian and Indian learning. Latin readers received a package wrapped in Arabic commentary. Knowledge was metabolised, not merely transported: each node digested, absorbed, and grew new tissue. Scholars call this \textbf{creative transformation}: translation as recreation.

Second, flow was multidirectional, not Greece towards Europe through an Arab rest stop. Numerals came from India, paper from China, medicine circulated among Persia, Syria, and Spain. What reached Europe was a worldwide composite. A network reduced to a one-way baton resembles a transport map reduced to one running track.

Third, custodians rejected faulty stock. Arabic astronomers did not simply accept Ptolemy; better observations revised his models and data across generations, with Samarkand tables far more accurate than ancient ones. Custodians do not correct customers' products; colleagues do. These scholars were Greeks' \textbf{peers and rivals}, not their custodians.

Should we therefore discard the relay and accept the network free of charge? No: dismantling an old framework does not exempt its replacement from scrutiny. Networks can hide unequal power under exchange and fusion: caliphs chose funding, conquerors' languages compelled change, Mongol swords removed nodes. Networks were never flat or equal. Warm stories of everyone's contribution can also conceal slavery and women's near-invisibility in surviving scholarly records. Good historical tools clarify rather than comfort.

One final layer: even the Islamic-world label is retrospective convenience. Christian Hunayn thought he was translating Galen, not participating in a project named Islamic civilisation. Samarkand papermakers thought about better pulp, not linking the Old World. We see the network from above; people inside follow individual routes. Civilisation is not one person's script but the aggregate of countless separate tasks.

\section[Algorithms, English, and Undersea Cables]{Back to Today: Algorithms, English, and Undersea Cables}
This thousand-year-old story repeats in your life, with you among its actors.

Your phone recommendations run algorithms, bearing al-Khwarizmi's name. Moving equation terms and changing signs enacts al-jabr, restoration. Your 1, 2, 3 were invented in India, modified by Arabs, and spread through Europe. Your paper followed China$\rightarrow$\allowbreak Samarkand$\rightarrow$\allowbreak Baghdad$\rightarrow$\allowbreak Cairo$\rightarrow$\allowbreak Spain$\rightarrow$\allowbreak the world. The chapter's entire cast now works in your schoolbag.

Consider English class. Shared scholarly Arabic was like today's scientific English, not merely metaphorically but as the same mechanism with a new part. Japanese and Brazilian chemists meet in English; Chinese laboratories publish in English journals; programming keywords are English. Today's network lingua franca plays Arabic's role. A young Persian studying Arabic and you memorising English a millennium later repeat the same act, including its frustration: Persian retreated from scholarship into poetry, as many languages now retreat from papers into daily life.

Geographical routes repeat too. Several Asia--Europe undersea internet cables follow old monsoon trade paths from Southeast Asia through the Indian Ocean, Red Sea, Suez, and Mediterranean. Wind-waiting ships and video-call fibres follow the same seabed. Light replaces wind and data replaces ledgers, while nodes remain Singapore, Mumbai, Aden, Suez, and Alexandria.

Schools offer another echo. China's recurring debates about translating more good foreign books and the relative value of translation and originality repeat House of Wisdom questions. They were never enemies. Al-Khwarizmi translated Greek and Indian astronomical tables and created original algebra. Translation connects nodes; originality grows them. Connection without growth shrinks nodes; growth without connection dissolves the net.

\section[Your Question: Bridge or Tollbooth?]{Your Question: Is a Common Language a Bridge or a Tollbooth?}
Return to the young Córdoban.

Born in Spain speaking a local dialect, he must learn Arabic for medical textbooks, teachers, examinations, practice, and reputation. He learns, joins the network, becomes renowned, and sees his works reach Cairo and Baghdad. He gains the world at a cost: his most important lifelong thoughts occur outside his mother tongue, which remains increasingly far beyond scholarship's doorway.

A millennium later it is your turn. English opens research and worldly windows, yet you may sense a barrier when things clearest in Chinese must become English before the world hears them.

Our final question is:

\textbf{When one language becomes knowledge's worldwide common language, does the network become fairer or less fair?}

For fairness, you have evidence: bright people anywhere need learn only one language to converse with humanity's leading minds. Admission costs a foreign language, far less than unaffordable parchment once did. Al-Khwarizmi from Khwarazm and Ibn Rushd from Córdoba travelled worldwide through that ticket.

For unfairness, you have evidence too. Tickets never cost nothing. Native speakers enter free; others pay years of youth. Their own languages retreat in knowledge meanwhile: most of today's languages can no longer produce a physics paper. Connecting everyone also quietly decides whose words count and which modes of thought appear professional.

Bridge or tollbooth? Answer with reasons. Look once more at your homework: paper, numerals, equations, three thousand-year travellers in your hands. You occupy a new node. Your generation will partly decide its next links, cargoes, and languages.
